% Standalone methodology section for inclusion in a larger LaTeX paper.
% Requires: \usepackage{amsmath,amssymb}

\section{Methodology}
\label{sec:methodology}

\subsection{Overview}

We propose \emph{BitSparse}, a GPU-resident activation-compression method for
feed-forward networks using ReLU or squared ReLU (ReLU$^2$) nonlinearities.
The method exploits the fact that ReLU activations are non-negative and often
contain many zeros.  Rather than retaining a full dense activation for the
backward pass, BitSparse stores a tile-wise support mask, a compact stream of
positive values, and a per-tile prefix index.  The representation preserves
the positions of active values while avoiding per-value coordinate indices.

BitSparse is designed primarily to reduce the memory required to save
activations during training.  It is not a sparse-GEMM method in its current
form: the forward feed-forward projection is computed densely, and the
sparse operand is reconstructed into a dense temporary before the
weight-gradient matrix multiplication.  The method therefore trades
compression and reconstruction work for reduced persistent activation memory.

\subsection{Tile-wise sparse activation format}
\label{sec:method-tile-format}

Let $X \in \mathbb{R}^{M \times N}$ denote a preactivation and let
$H = \operatorname{ReLU}(X)$.  BitSparse partitions $H$ into a regular grid
of tiles of size $B_M \times B_N$, where the implementation uses
$B_M=B_N=64$.  The number of tile rows and columns is
\begin{equation}
    G_M = \left\lceil \frac{M}{B_M} \right\rceil,
    \qquad
    G_N = \left\lceil \frac{N}{B_N} \right\rceil.
\end{equation}
Tiles are ordered in row-major grid order, such that the identifier of tile
$(u,v)$ is $t=uG_N+v$.  Boundary tiles are logically padded with zeros;
therefore, every tile has $B_MB_N=4096$ positions irrespective of its
location.

For each position $q$ in tile $t$, we form a binary support indicator
\begin{equation}
    m_t(q) = \mathbb{I}[H_t(q)>0].
\end{equation}
The representation of $H$ consists of the following arrays:
\begin{enumerate}
    \item a packed bitmask array $\mathbf{m}$, containing one bit per
    logical activation position;
    \item a compact value array $\mathbf{v}$, containing only positive
    activation values in tile-major and row-major-within-tile order; and
    \item a 32-bit prefix array $\mathbf{p}$ of length $G_MG_N+1$.
\end{enumerate}
The prefix array is defined by
\begin{equation}
    p_0=0,
    \qquad
    p_{t+1}=p_t+\sum_{q=0}^{B_MB_N-1}m_t(q).
    \label{eq:prefix}
\end{equation}
Thus, $p_t$ is the starting logical offset of tile $t$ in the compact value
stream.  The rank of an active local position is
\begin{equation}
    r_t(q)=\sum_{j=0}^{q-1}m_t(j),
\end{equation}
and its stored value is
\begin{equation}
    v_{p_t+r_t(q)} = H_t(q)
    \qquad \text{when } m_t(q)=1.
    \label{eq:value-address}
\end{equation}
Equations~\eqref{eq:prefix} and~\eqref{eq:value-address} permit direct
addressing within a tile without storing an explicit row and column index for
each nonzero.

The bitmask is stored as unsigned bytes.  If $f$ is the flattened position in
a tile, its indicator is held in bit $f \bmod 8$ of byte
$\lfloor f/8 \rfloor$.  Consequently, each $64\times64$ tile requires
512 bytes of mask metadata.  The full sparse object additionally records the
logical shape, tile geometry, stored dtype, optional FP8 scale, and an offset
into a shared value buffer when one is used.

\subsection{GPU compression and reconstruction}
\label{sec:method-kernels}

Compression is performed entirely on the GPU in three stages.  First, a
Triton tile kernel loads one tile, evaluates $H_t>0$, packs groups of eight
Boolean values into mask bytes, and reduces the same mask to a nonzero count.
Second, an inclusive scan over the per-tile counts constructs the prefix
array.  The implementation initializes the first prefix element on the
device and writes the scan through a signed 32-bit view of the unsigned
prefix storage; this avoids a synchronizing scalar assignment from the host.
Third, a second Triton kernel expands each tile mask, computes
$r_t(q)$ using a tile-local prefix scan, and scatters active values into the
compact stream according to Equation~\eqref{eq:value-address}.

To reconstruct a tile, the inverse kernel expands the packed mask, recomputes
the local ranks, gathers the corresponding compact values, and stores zeros
at inactive positions.  The same kernel supports a batched reconstruction of
contiguous tile rows.  This enables a memory-limited alternative in which
only a row-aligned slice of a sparse matrix is reconstructed at one time.

\subsection{Sign-bit packing}
\label{sec:method-packing}

Because stored ReLU values are non-negative, their sign bit is always zero.
BitSparse can therefore remove this redundant bit and concatenate the
remaining bits into a continuous stream.  BF16 values occupy 15 bits rather
than 16 bits, while FP8 values occupy 7 bits rather than 8 bits.  This
operation is bit-exact with respect to the already selected BF16 or FP8
storage format; it is distinct from numerical quantization.

Let $b$ denote the number of stored bits per value ($b=15$ for BF16 and
$b=7$ for FP8).  The bit position of compact value $i$ is $ib$.  During
decoding, the implementation reads two adjacent 16-bit words and computes
\begin{equation}
    w = \left\lfloor \frac{ib}{16} \right\rfloor,
    \qquad
    s = ib \bmod 16,
    \qquad
    \hat{v}_i =
    \left((W_w \mathbin{|} (W_{w+1} \ll 16)) \gg s\right)
    \mathbin{\&} (2^b-1),
\end{equation}
followed by a bit reinterpretation into the appropriate BF16 or FP8 type.
Two-word reads handle values that cross word boundaries.  Standalone packed
arrays include a guard word, while value streams shared across layers align
each activation to eight logical values.  This alignment makes every
independent packed stream byte-aligned and prevents neighboring pack kernels
from updating the same output byte.

\subsection{FP8 activation storage and matrix multiplication}
\label{sec:method-fp8}

BitSparse supports BF16 and FP8 activation storage.  For FP8 conversion, a
single tensor-wise scale is computed as
\begin{equation}
    s = \max\left(
        \frac{\max_{i,j}|H_{ij}|}{q_{\max}},
        10^{-9}
    \right),
    \qquad
    H_{\mathrm{FP8}} =
    \operatorname{cast}_{\mathrm{FP8}}\left(\frac{H}{s}\right),
    \label{eq:fp8-scale}
\end{equation}
where $q_{\max}$ is the largest finite value of the selected FP8 format.
The implementation currently converts tensors to E4M3FN FP8.  The scale is
stored as sparse metadata and is applied by the matrix multiplication or
ReLU$^2$ reconstruction path as appropriate.

Dense FP8 products use PyTorch's scaled matrix multiplication primitive with
tensor-wise scales.  The implementation ensures that the left operand is
row-major contiguous and applies an architecture-dependent layout conversion
to the right operand when required by cuBLASLt.  Specifically, the right
operand is reformatted to a column-major-compatible layout on Ada-class
devices, while this copy is skipped on Blackwell-class devices that accept
the row-major layout directly.  FP8 matrix multiplication returns BF16
outputs.

\subsection{Sparse activation checkpointing for feed-forward layers}
\label{sec:method-ffn}

We integrate the representation into two custom PyTorch autograd functions.
For an input $Z$, first-layer weight
$W_1 \in \mathbb{R}^{d_{\mathrm{ff}}\times d_{\mathrm{in}}}$, and
second-layer weight
$W_2 \in \mathbb{R}^{d_{\mathrm{out}}\times d_{\mathrm{ff}}}$, the
ReLU feed-forward block computes
\begin{align}
    A &= ZW_1^\top + b_1, \\
    H &= \operatorname{ReLU}(A), \\
    Y &= HW_2^\top + b_2.
\end{align}
The dense $H$ is used by the forward projection, but its BitSparse
representation is saved for backward in place of a full dense saved
activation.

For the ReLU$^2$ variant, the saved sparse tensor still represents
$H=\operatorname{ReLU}(A)$, whereas the projected activation is
\begin{equation}
    Y = \left(\kappa H^2\right) W_2^\top + b_2,
    \label{eq:relu2-forward}
\end{equation}
where $\kappa$ is \texttt{RELU2\_SCALE} (equal to one in the current
implementation).  In the FP8 ReLU$^2$ path, the compressed activation is an
FP8 copy used for backward, while the forward square is formed from the
original dense ReLU activation before the output matmul.  Consequently, FP8
backward computations use a quantized approximation of $H$, whereas sign-bit
packing introduces no additional numerical approximation.

Let $G=\partial \mathcal{L}/\partial Y$ denote the upstream gradient.  For
the ReLU block, the second-layer weight gradient is
\begin{equation}
    \frac{\partial \mathcal{L}}{\partial W_2}
    = G^\top H.
\end{equation}
The implementation reconstructs $H$ from its BitSparse form and then
performs this dense or scaled-FP8 product.  The preactivation gradient is
computed as
\begin{equation}
    \frac{\partial \mathcal{L}}{\partial A}
    = (GW_2) \odot \mathbb{I}[H>0].
    \label{eq:relu-backward}
\end{equation}
Rather than reconstructing $H$ to apply the indicator in
Equation~\eqref{eq:relu-backward}, a Triton kernel expands the stored
bitmask and overwrites inactive gradient entries with zero.

For the ReLU$^2$ block, the corresponding gradients are
\begin{align}
    \frac{\partial \mathcal{L}}{\partial W_2}
    &= G^\top (\kappa H^2), \\
    \frac{\partial \mathcal{L}}{\partial A}
    &= (GW_2) \odot 2\kappa H \odot \mathbb{I}[H>0].
    \label{eq:relu2-backward}
\end{align}
The second-layer weight gradient reconstructs $\kappa H^2$ in BF16, which
avoids overflow from squaring FP8 values.  For the second expression in
Equation~\eqref{eq:relu2-backward}, a fused Triton kernel reads the compact
values directly, applies $2\kappa H$ at active positions, and writes zero
at inactive positions.  Gradients with respect to $W_1$ and $Z$ are then
computed by PyTorch from the returned preactivation gradient.

\subsection{Shared value buffers and implementation details}
\label{sec:method-buffer}

When processing multiple feed-forward layers, separate compact value arrays
would create substantial allocation overhead.  BitSparse optionally appends
the value streams of multiple activations to a preallocated
\texttt{TensorBuffer}.  The buffer holds a device-resident logical offset;
each sparse tensor records a copy of its starting offset, and compression
advances the shared offset by its active-value count.  Packed streams are
rounded up to the next eight-value boundary.  The buffer is reset only after
all forward activations using it have completed backward, since resetting it
earlier would overwrite activations referenced by the autograd graph.

All tile packing, compaction, unpacking, and gradient masking kernels are
implemented in Triton.  Kernel configurations are autotuned over small sets
of warp counts and pipeline stages.  The unpack and gradient kernels operate
on one tile per program, whereas the packing kernel is tuned over block size,
warp count, and storage codec.  The implementation uses in-place operations
for ReLU, ReLU$^2$ scaling, and gradient masking where safe, further reducing
temporary allocation pressure.
